\documentclass[11pt]{amsart}
\usepackage[margin=1in]{geometry}
\usepackage{amsmath,amssymb,mathtools,hyperref}
\newtheorem{theorem}{Theorem}
\newtheorem{lemma}[theorem]{Lemma}
\newtheorem{corollary}[theorem]{Corollary}
\newtheorem{proposition}[theorem]{Proposition}
\theoremstyle{definition}\newtheorem{definition}[theorem]{Definition}
\newcommand{\C}{\mathbb C}
\newcommand{\one}{\mathbf1}
\newcommand{\VEC}{\operatorname{Vec}}
\newcommand{\Hom}{\operatorname{Hom}}
\newcommand{\rad}{\operatorname{rad}}
\newcommand{\id}{\operatorname{id}}
\newcommand{\Pauli}{\mathcal P}
\title{Adaptive charge measurements in abelian twisted finite gauge theories}
\author{Titania}
\date{14 November 2026. Candidate for whole assessment; not submitted.}
\begin{document}
\begin{abstract}
We give a proposed classical simulation theorem for intrinsic planar anyon
protocols in every twisted double of a finite abelian group. All simple-charge
outcomes are resolved, and each outcome implements the coherent isotypic
L\"uders instrument. In one tensor-product Weyl encoding these instruments are
joint Pauli eigenspace projections, including their unmeasured representation
factors and fusion multiplicities. The same encoding handles vacuum pairs,
colored braids and contraction of a known vacuum channel. This would also cover
untwisted doubles of every finite group of nilpotency class at most two, by a
prior pointed Morita equivalence. An independent binary vacuum/nonvacuum test
is a different resource and is excluded. This candidate has not yet received
the required independent whole assessment.
\end{abstract}
\maketitle

\section{Protocol and statement}
Fix a unitary modular category. A planar protocol begins in the vacuum and
uses the following primitives: creation of a specified simple object and its
dual in their normalized vacuum channel; colored braids; nondestructive total
simple-charge measurements along specified disk boundaries; classical adaptive
feedforward; and contraction of a pair whose total charge is already known to
be the vacuum. Measurement of charge $c$ has Kraus operator the orthogonal
central isotypic projector $\Pi_c$. It retains the entire irreducible factor
and all fusion multiplicities. There is no measurement of a finer fusion basis.
Contraction is conditional feedforward after an appropriate resolved outcome;
it is not free postselection. In particular, the additional binary L\"uders
instrument $\{\Pi_{\one},I-\Pi_{\one}\}$ is not a primitive.

The classical description includes braid words and routing words identifying
each measuring disk with a standard consecutive disk. Its total length is $L$;
the maximum number of elementary anyons is $n$. Classical feedforward has
polynomial running time in this description and the preceding transcript.
Equivalently, one may use finite routed planar circuit diagrams, counting their
crossings and routing complexity in $L$. Local types and the category are fixed.

\begin{theorem}\label{thm:main}
Let $A$ be any finite abelian group and let $[\omega]\in H^3(A;U(1))$.
For the unitary category $Z(\VEC_A^\omega)$, the protocols just defined admit
exact classical sampling with expected running time polynomial in $n,L$.
They admit sampling to total-variation error at most $\epsilon$ in bounded
time polynomial in $n,L,\log(1/\epsilon)$. Constants and preprocessing may
depend on $A,[\omega]$. Along every nonzero transcript, the pure conditional
state has one efficiently updated generalized Pauli stabilizer description.
\end{theorem}

\begin{corollary}\label{cor:nil2}
The same conclusion holds for $Z(\VEC_G)$ for every fixed finite group $G$
with $[G,G]\subset Z(G)$. In particular it covers all groups of class at most
two and exponent an odd prime, including every $UT_3(\mathbb F_p)$.
\end{corollary}

The finite braid-image theorem of Etingof--Rowell--Witherspoon~\cite{ERW}
does not supply this conclusion: measurements change the available resources.
Mochon's universality theorem for solvable nonnilpotent finite groups has specified
calibrated ancillary resources~\cite{Mochon}. Galindo--Rowell--Wang~\cite{GRW}
identify acyclic and nilpotent fusion rules and ask about measurement-assisted
computational power. The classical adaptive normalizer simulation engine is
prior work of Bermejo-Vega--Van den Nest~\cite{BV}.
Ground-state preparation~\cite{TVV}, Clifford-hierarchy lattice realizations
and code-switching magic~\cite{Fechisin}, and selected quasistabilizer Lindblad
trajectories~\cite{Chirame} are different instrument statements. Our proposed
new content is their missing intrinsic all-simple-isotypic compatibility theorem.
The underlying Weyl representation theory and categorical equivalences are prior.

\section{A simultaneous Weyl model}
For a finite abelian group $K$, put $\mathcal H_K=\C[K]$, with orthonormal
basis $|k\rangle$. Define $X_u|k\rangle=|k+u\rangle$ and
$Z_\chi|k\rangle=\chi(k)|k\rangle$ for $\chi\in\widehat K$.
We call scalar multiples of $X_uZ_\chi$ projective Paulis. Scalars never change
an eigenspace measurement, but their exact values are retained in the labels.
A unitary normalizing this projective Pauli group is Clifford. For the
simulation we specify a Clifford by its conjugation on Pauli generators.

\begin{lemma}\label{lem:heis}
Let $A$ be finite abelian and $\beta\in Z^2(A;U(1))$. Set
$B(x,y)=\beta(x,y)/\beta(y,x)$ and $R=\rad B$.
There are a finite abelian group $K$, a projective Pauli representation
$W:A\to U(\mathcal H_K)$ whose kernel modulo scalars is $R$, and a function
$q:A\to U(1)$ such that all irreducible $\beta$-projective representations are
\[
 \pi_\lambda(x)=q(x)\widetilde\lambda(x)W(x),\qquad
 \lambda\in\widehat R,
\]
where $\widetilde\lambda$ is any character extension to $A$.
We may take $W(r)=I$ for $r\in R$. Every irreducible has dimension
$|K|=\sqrt{|A/R|}$. In the twisted group algebra the elements
$q(r)^{-1}e_r$, $r\in R$, form its central group algebra $\C[R]$.
\end{lemma}
\begin{proof}
The commutator $B$ is an alternating bicharacter. Its quotient on $A/R$ is
nondegenerate. A finite abelian group with a nondegenerate alternating
bicharacter is an orthogonal sum of hyperbolic cyclic pairs: choose an element
of maximum order, pair it with a character partner of the same order, split
off the resulting nondegenerate pair, and induct. Thus
$A/R\cong K\oplus\widehat K$ with its standard pairing.
Pulling back $X_uZ_\chi$ gives $W$ with a multiplier $c$ having commutator $B$.
The symmetric cocycle $\beta/c$ is a coboundary. Indeed its abelian central
extension by $U(1)$ splits because $U(1)$ is divisible and therefore injective
as an abelian group. Choose $q$ with $\delta q=\beta/c$.
Characters of $R$ extend to $A$ by the same injectivity.

The centre of $\C^\beta[A]$ is spanned by the $e_r$, $r\in R$, since commuting
with $e_x$ multiplies a basis element by $B(x,r)$. The normalized elements
$q(r)^{-1}e_r$ multiply as $R$. Fixing their character $\lambda$ leaves an
algebra of dimension $|A/R|$ with scalar centre. As a finite-dimensional
semisimple complex algebra it is $\operatorname{Mat}_{|K|}(\C)$.
The Weyl model displayed above realizes this block and all its characters.
All matrices are unitary, and the image of $A$ modulo scalars is the entire
projective Pauli group on $\mathcal H_K$.
\end{proof}

We use the quasi-Hopf conventions in the published \cite[Section 3]{ERW}.
For additive abelian $A$, define
\begin{align*}
 \beta_a(x,y)&=\frac{\omega(a,x,y)\omega(x,y,a)}{\omega(x,a,y)},\\
 \gamma_x(a,b)&=\frac{\omega(a,b,x)\omega(x,a,b)}{\omega(a,x,b)}.
\end{align*}
A simple is a fixed flux $a$ and an irreducible $\beta_a$-projective charge.
The tensor charge is
\begin{equation}\label{eq:coproduct}
 \rho_{X\otimes Y}(x)=\gamma_x(a,b)\rho_X(x)\otimes\rho_Y(x).
\end{equation}
The associator on fixed fluxes is a scalar of modulus one, and
\begin{equation}\label{eq:braid}
 c_{X,Y}=\operatorname{flip}\circ(I\otimes\rho_Y(a)).
\end{equation}
These are prior structural formulas, also recorded in inverse-cocycle
conventions by Bantay~\cite[Appendix A]{Bantay}.
For each flux $a$, make the choices of Lemma~\ref{lem:heis} once, writing
$B_a,R_a,K_a,W_a,q_a$. Elementary-anyon registers are $\C[K_a]$.

\section{The exact coarse instrument}
\begin{lemma}\label{lem:projector}
Let $X_1,\ldots,X_m$ be fixed simples with fluxes $a_i$, and put
$a=\sum_i a_i$. On the entire tensor register
$\mathcal H=\bigotimes_i\C[K_{a_i}]$, let $\rho(x)$ be its actual charge action,
including every factor in \eqref{eq:coproduct}. For $r\in R_a$ put
$U(r)=q_a(r)^{-1}\rho(r)$. Then $U$ is an ordinary representation of $R_a$
by commuting projective Paulis. The full isotypic projector for charge
$\lambda\in\widehat R_a$ is
\begin{equation}\label{eq:projector}
 \Pi_{a,\lambda}=\frac1{|R_a|}\sum_{r\in R_a}\overline{\lambda(r)}U(r).
\end{equation}
Measuring a generating set of these commuting operators implements exactly
the L\"uders instrument $v\mapsto\Pi_{a,\lambda}v$, with all multiplicities intact.
\end{lemma}
\begin{proof}
Each $\rho_{X_i}(x)$ is a scalar Pauli by Lemma~\ref{lem:heis}. The tensor
corrections in \eqref{eq:coproduct} are scalars, so the same holds for $\rho(x)$.
Since the coproduct is an algebra homomorphism, its multiplier is precisely
$\beta_a$. Its restriction to the central normalized units
$q_a(r)^{-1}e_r$ is therefore the ordinary commuting representation $U$.
Lemma~\ref{lem:heis} identifies these central characters with all simple
isotypes, proving \eqref{eq:projector}. This is a central idempotent, hence is
the identity on the entire irreducible factor and multiplicity of that isotype.
All operators are unitary, so the idempotents are orthogonal projectors.

The tuple of eigenvalues on generators of $R_a$ determines exactly one
character of $R_a$, hence exactly one simple charge. Sequentially measuring
these commuting generators has joint Kraus projector \eqref{eq:projector}.
It does not distinguish vectors within an isotype. When $U(r)$ is a scalar
times a standard Pauli, measure that standard Pauli and restore the known
scalar in the classical outcome label. This changes no eigenspace. Redundant
or deterministic generators add no information. In particular, the argument
does not refine a charge into irreducible-vector or multiplicity outcomes.
\end{proof}

For total flux zero choose $q_0=1$. The physical space
$\Hom(\one,\bigotimes_i X_i)$ is the invariant subspace for $\rho(x)$,
$x\in A$, and is therefore a joint Pauli eigenspace, a stabilizer code.
A standard subdisk's central idempotent commutes with the full algebra action,
so preserves this physical space. Alternatively its Pauli commutator with a
global charge is $B_a(r,x)=1$ for $r\in R_a$.
Reassociation is scalar on each fixed elementary-flux word.
Equation~\eqref{eq:braid} is a register swap followed by a Pauli on one factor.
It is Clifford and carries the physical invariant space to the correctly
reordered one by the quasitriangular identity. A routed disk instrument is
therefore a Clifford conjugate of Lemma~\ref{lem:projector}, in the same
persistent tensor registers. No fusion-basis change or new polarization is
needed after a measurement.

\section{Dual cups and caps, including 2-torsion}
\begin{lemma}\label{lem:cup}
In the choices of elementary registers above, every normalized simple/dual
vacuum cup is a Clifford image of a maximally entangled pair. A cap on an
already known vacuum channel is the inverse stabilizer isometry.
\end{lemma}
\begin{proof}
If $X$ has flux $a$ and charge $\pi$, then on the conjugate Hilbert space set
\[
 \pi^\vee(x)=\gamma_x(a,-a)^{-1}\overline{\pi(x)}.
\]
The identity asserting that \eqref{eq:coproduct} has multiplier $\beta_0=1$
shows that this has multiplier $\beta_{-a}$. The normalized Bell vector is
invariant for $\gamma_x(a,-a)\pi(x)\otimes\pi^\vee(x)$. It consequently
spans the unique vacuum channel and realizes the dual object. The scalar
rigidity corrections only change the phase of its normalized cup.

The commutator form for $-a$ is the inverse of that for $a$, so choose the
same abstract $K$ for the two elementary registers. The conjugates of standard
Paulis are standard Paulis. The desired representation $\pi^\vee$ and the
already chosen irreducible representation of its simple type are unitarily
intertwined. Let $V$ be an intertwiner. Both representations projectively
surject onto the Pauli group, by Lemma~\ref{lem:heis}. Hence conjugation by $V$
carries the entire projective Pauli group onto itself: $V$ is Clifford.
This reasoning also applies when $a=-a$ and when the simple is self-dual;
it does not require a separate basis choice for the same object. Orthogonal
or symplectic self-duality changes the intertwiner but not its Clifford property.
The actual cup is thus $(I\otimes V)$ applied to Bell, up to phase and choice
of direction for $V$.

Bell on $\C[K]\otimes\C[K]$ is stabilized by $X_u\otimes X_u$ and
$Z_\chi\otimes Z_{\bar\chi}$. Its preparation and inverse on its known
one-dimensional channel are stabilizer operations. A known Bell factor is
unentangled from the complementary registers; applying the inverse Clifford
and discarding that factor implements the normalized cap.
\end{proof}

\section{Compilation and the group corollary}
\begin{proof}[Proof of Theorem~\ref{thm:main}]
Precompute the finite charge models, cocycles, dual intertwiners and their
Pauli conjugation data. Choose a root-of-unity representative of $[\omega]$
using \cite[Lemma 4.3]{ERW}. The symmetric-coboundary equations defining
$q_a$ are a finite integer linear system for phases with rational right-hand
side, so Smith normal form supplies root-of-unity choices. Thus every scalar
table is encoded by integer exponents modulo a fixed denominator. Conjugation
data for a dual intertwiner are obtained by lifting each Pauli generator through
the surjection $A/R_a$ in Lemma~\ref{lem:heis} and using its intertwining identity.
All these constants depend only on the fixed category.
A tensor word uses $\C[K]$ for $K=\prod_i K_{a_i}$, with
$\log|K|=O(n)$. Start with the vacuum stabilizer state. By Lemma~\ref{lem:cup},
each allowed pair creation adds a stabilizer Bell pair followed by a Clifford.
Colored braids and reassociations preserve the Pauli description.
Lemma~\ref{lem:projector} realizes exactly every allowed charge instrument,
including arbitrary later routed disks. Known-vacuum contraction is covered
by Lemma~\ref{lem:cup}. Induction gives one stabilizer description on every
nonzero transcript, with all physical multiplicities retained.

The generalized stabilizer update and sampling algorithms of \cite{BV}
therefore apply. For unitary updates their proof needs computable conjugation
on Pauli generators; here it is explicit for braids and finite precomputed
dual intertwiners. It does not require searching for a decomposition of a
Clifford into gates. Commuting projective Pauli measurements are standard
Pauli measurements with known scalar relabeling, as in the projector lemma.
Their number per disk is bounded by the fixed group $A$.
Thus their stabilizer algorithms have polynomial cost in $n,L$.
Adaptive feedforward chooses the next operation from the sampled transcript.
The only unbounded runtime is exact finite-uniform random sampling, implementable
by rejection from random bits. Its exponential rejection tail gives the stated
bounded-time approximation after $O(\log(L/\epsilon))$ attempts per sample.
No conditioning on a rare successful transcript has been granted for free.
\end{proof}

\begin{proof}[Proof of Corollary~\ref{cor:nil2}]
Put $Z=Z(G)$ and $Q=G/Z$. The standard pointed Morita equivalence using
the subgroup $Z$ identifies the dual of $\VEC_G$ with
$\VEC_{\widehat Z\times Q}^{\eta}$, where $\eta$ is a unitary associator.
This is the central-subgroup, trivial-cochain specialization of
\cite[Theorems 4.5 and 4.9]{Naidu}; in a section of $G\to Q$ with central
extension cocycle $\kappa$, its associator is
$\eta((\chi,x),(\psi,y),(\xi,z))=\chi(\kappa(y,z))$.
Since $G$ has class at most two, $Q$ and $\widehat Z\times Q$ are abelian.
Taking centres gives the braided equivalence used here, also recorded in
\cite[Theorem D.1]{Fechisin}. Both categories are unitary; Reutter's
\cite[Theorem 5]{Reutter} makes the braided equivalence unitary.
It therefore preserves inner products, normalized cups, braids and orthogonal
simple-isotypic projectors, hence the exact protocol and its probabilities.
All relabeling and local equivalence data are fixed-category preprocessing.
Apply Theorem~\ref{thm:main}.
\end{proof}

\section{Controls and the additional binary resource}
For $A=(\mathbb Z/4)^3$ and $\omega(a,b,c)=i^{a_1b_2c_3}$, take two
flux-$e_1$ charge spaces with $\pi(x)=X^{x_3}Z^{x_2}$ on $\C[\mathbb Z/4]$.
Their total flux is $2e_1$. Its radical is generated by $e_1,2e_2,2e_3$;
the first acts trivially, and the other two act as $Z^2\otimes Z^2$ and
$X^2\otimes X^2$. There are four nonzero coarse isotypic projections of rank
$4$, each retaining an irreducible factor of dimension $2$ and multiplicity $2$.
This tests a changing Weyl-register size and a genuinely unmeasured multiplicity.
For $A=\mathbb F_2^3$, the same type-III formula and flux $(1,1,1)$ give
a two-dimensional charge with a self-dual cup whose intertwiner is $XZ$.
It is antisymmetric but Clifford, illustrating that a self-dual 2-torsion
sector is included in Lemma~\ref{lem:cup}.

The restriction to resolved simple-charge outcomes has mathematical content.
For the $UT_3(\mathbb F_p)$ model take $A=\mathbb F_p^3$ and
$\omega(a,b,c)=\exp(2\pi i a_1b_2c_3/p)$, $p$ odd.
At flux $a=(1,0,0)$ an irreducible charge has dimension $p$.
Choose dual coordinates so cups are Bell and form two vacuum pairs with
colors $a,-a,a,-a$:
\[
 \psi=\frac1p\sum_{i,j}|i,i,j,j\rangle.
\]
The middle pair's vacuum projector is Bell projection. Consequently
\[
 \Pi_{\one,23}\psi=\frac1{p^2}\sum_{i,k}|i,k,k,i\rangle.
\]
The failed binary branch $(I-\Pi_{\one,23})\psi$ has support size
$2p^2-p$ and probability $1-p^{-2}$. Its support size is not a power of $p$,
so this vector is not a prime-qudit Pauli stabilizer state. All vectors lie
in the physical total-vacuum sector. This is an actual additional-instrument
failure of the proposed Pauli closure, not a claim of universality or of
classical hardness. It does not occur in the resolved-instrument theorem.

\begin{thebibliography}{99}
\bibitem{ERW} P. Etingof, E. Rowell, S. Witherspoon, Braid group representations
from twisted quantum doubles of finite groups, Pacific J. Math. 234 (2008), 33--41.
\url{https://msp.org/pjm/2008/234-1/pjm-v234-n1-p03-s.pdf}.
\bibitem{Bantay} P. Bantay, Algebraic aspects of orbifold models,
\url{https://arxiv.org/abs/hep-th/9303009}.
\bibitem{BV} J. Bermejo-Vega, M. Van den Nest, Classical simulations of
Abelian-group normalizer circuits with intermediate measurements,
Quantum Inf. Comput. 14 (2014), 181--216.
\url{https://arxiv.org/abs/1210.3637}.
\bibitem{Mochon} C. Mochon, Anyon computers with smaller groups,
Phys. Rev. A 69 (2004), 032306. \url{https://arxiv.org/abs/quant-ph/0306063}.
\bibitem{GRW} C. Galindo, E. Rowell, Z. Wang, Acyclic anyon models,
thermal anyon error corrections, and braiding universality,
\url{https://arxiv.org/abs/1801.04296}.
\bibitem{TVV} N. Tantivasadakarn, A. Vishwanath, R. Verresen,
Hierarchy of topological order from finite-depth unitaries, measurement,
and feedforward, PRX Quantum 4 (2023), 020339.
\url{https://arxiv.org/abs/2209.06202}.
\bibitem{Fechisin} C. Fechisin et al., Clifford-hierarchy stabilizer formalism
with applications to twisted quantum doubles,
\url{https://arxiv.org/abs/2610.00464}.
\bibitem{Chirame} S. Chirame, A. Prem, S. Gopalakrishnan, F. Burnell,
Stabilizing non-Abelian topological order against heralded noise via local
Lindbladian dynamics, PRX Quantum 6 (2025), 030363.
\url{https://arxiv.org/abs/2410.21402}.
\bibitem{Naidu} D. Naidu, Categorical Morita equivalence for group-theoretical
categories, Comm. Algebra 35 (2007), 3544--3565.
\url{https://arxiv.org/abs/math/0605530}.
\bibitem{Reutter} D. Reutter, Uniqueness of unitary structure for unitarizable
fusion categories, Comm. Math. Phys. (2023).
\url{https://arxiv.org/abs/1906.09710}.
\end{thebibliography}
\end{document}
